% Preámbulo tipográfico del libro Medicina y Ciencia (título y autoría: datos-libro.tex).
% Maqueta de catálogo KDP / IngramSpark: 6,69 × 9,61 pulgadas = 170 × 244 mm (scrbook, dos caras).
% Compilado con LuaLaTeX a través de Quarto.

\usepackage{microtype}
\usepackage{booktabs}
\usepackage{longtable}
\usepackage{array}
\usepackage{xcolor}
\usepackage{scrlayer-scrpage}
\usepackage{enumitem}
\usepackage{etoolbox}
\usepackage{needspace}
\usepackage{tikz}
\usepackage[most]{tcolorbox}

% Tipografía secundaria para cabeceras y rótulos.
\newfontfamily\displayfont{Libertinus Serif Display}[
  NFSSFamily=libdisplay,
  Ligatures=TeX
]
\newfontfamily\fallbackserif{FreeSerif}[
  Scale=MatchLowercase
]

% Símbolos Unicode clínicos frecuentes.
\usepackage{newunicodechar}
\newunicodechar{≥}{\ensuremath{\geq}}
\newunicodechar{≤}{\ensuremath{\leq}}
\newunicodechar{≠}{\ensuremath{\neq}}
\newunicodechar{≈}{\ensuremath{\approx}}
\newunicodechar{±}{\ensuremath{\pm}}
\newunicodechar{×}{\ensuremath{\times}}
\newunicodechar{−}{\ensuremath{-}}
\newunicodechar{→}{\ensuremath{\rightarrow}}
\newunicodechar{←}{\ensuremath{\leftarrow}}
\newunicodechar{↔}{\ensuremath{\leftrightarrow}}
\newunicodechar{↑}{\ensuremath{\uparrow}}
\newunicodechar{↓}{\ensuremath{\downarrow}}
\newunicodechar{α}{\ensuremath{\alpha}}
\newunicodechar{β}{\ensuremath{\beta}}
\newunicodechar{γ}{\ensuremath{\gamma}}
\newunicodechar{δ}{\ensuremath{\delta}}
\newunicodechar{μ}{\ensuremath{\mu}}
\newunicodechar{▸}{{\fallbackserif ▸}}
\newunicodechar{✓}{{\fallbackserif ✓}}
\newunicodechar{✗}{{\fallbackserif ✗}}

% Paleta editorial PowerSemiotics — Medicina y Ciencia (azul noche, oro cálido y borgoña clínico).
\definecolor{libroprimario}{HTML}{14283D}
\definecolor{libroacento}{HTML}{7A2E1D}
\definecolor{librogris}{HTML}{4A5568}
\definecolor{librofondo}{HTML}{F7F8FA}
\definecolor{libroborde}{HTML}{CBD5E1}
\definecolor{libroalerta}{HTML}{FDF5F3}
\definecolor{librosocratico}{HTML}{F2F6FA}
\definecolor{librosocraticoborde}{HTML}{2B547E}
\definecolor{portadafondo}{HTML}{0D1B2A}
\definecolor{portadapapel}{HTML}{F8F6F1}
\definecolor{portadaoro}{HTML}{C69C6D}
\definecolor{portadacian}{HTML}{4EA8DE}

% Interlineado y ritmo de párrafo de libro impreso (170 × 244 mm).
\linespread{1.11}
\setlength{\parindent}{1.15em}
\setlength{\parskip}{0pt}
\raggedbottom

% Control de viudas, huérfanas y cortes de página.
\widowpenalty=10000
\clubpenalty=10000
\displaywidowpenalty=10000
\predisplaypenalty=10000
\postdisplaypenalty=5000
\brokenpenalty=5000
\emergencystretch=2em
\hyphenpenalty=800
\exhyphenpenalty=400

% Listas compactas con buen aire vertical.
\setlist[itemize]{
  topsep=0.35\baselineskip,
  itemsep=0.18\baselineskip,
  parsep=0pt,
  leftmargin=1.4em
}
\setlist[enumerate]{
  topsep=0.35\baselineskip,
  itemsep=0.20\baselineskip,
  parsep=0pt,
  leftmargin=1.6em
}
\setlist[description]{
  topsep=0.4\baselineskip,
  itemsep=0.25\baselineskip,
  font=\normalfont\sffamily\bfseries\color{libroprimario}
}

% Jerarquía KOMA-Script.
\addtokomafont{disposition}{\color{libroprimario}}
\setkomafont{part}{\displayfont\Huge\bfseries\color{libroprimario}}
\setkomafont{partnumber}{\sffamily\normalsize\mdseries\scshape\color{libroacento}}
\setkomafont{chapter}{\displayfont\LARGE\bfseries\color{libroprimario}}
\setkomafont{chapterprefix}{\sffamily\footnotesize\mdseries\scshape\color{libroacento}}
\setkomafont{section}{\normalfont\large\bfseries\color{libroprimario}}
\setkomafont{subsection}{\normalfont\normalsize\bfseries\color{libroprimario}}
\setkomafont{subsubsection}{\normalfont\small\sffamily\bfseries\color{librogris}}
\setkomafont{pageheadfoot}{\normalfont\footnotesize\scshape\color{librogris}}
\setkomafont{pagenumber}{\normalfont\footnotesize\color{libroprimario}}

% Formato de capítulo con filete fino bajo el título.
\KOMAoptions{chapterprefix=true}
\renewcommand*{\chapterformat}{%
  \mbox{\chapappifchapterprefix{\nobreakspace}\thechapter\autodot\IfUsePrefixLine{}{\enskip}}%
}
\renewcommand*{\chapterlineswithprefixformat}[3]{%
  \raggedright
  {\usekomafont{chapterprefix}#2\par}%
  \vspace{0.2\baselineskip}%
  {\usekomafont{chapter}#3\par}%
  \vspace{0.45\baselineskip}%
  {\color{libroacento}\rule{\linewidth}{0.6pt}\par}%
}

\RedeclareSectionCommand[
  beforeskip=-1.8\baselineskip,
  afterskip=1.4\baselineskip
]{chapter}
\RedeclareSectionCommand[
  beforeskip=-1.3\baselineskip,
  afterskip=0.55\baselineskip
]{section}
\RedeclareSectionCommand[
  beforeskip=-0.95\baselineskip,
  afterskip=0.35\baselineskip
]{subsection}
\RedeclareSectionCommand[
  beforeskip=-0.7\baselineskip,
  afterskip=0.25\baselineskip
]{subsubsection}

% Evitar títulos sueltos al final de la página.
\preto\section{\needspace{5\baselineskip}}
\preto\subsection{\needspace{4\baselineskip}}
\preto\subsubsection{\needspace{3\baselineskip}}

% Cabeceras y pies de página (scrbook twoside).
\clearpairofpagestyles
\lehead{\headmark}
\rohead{\headmark}
\lefoot[\pagemark]{\pagemark}
\rofoot[\pagemark]{\pagemark}
\automark[section]{chapter}
\pagestyle{scrheadings}

% Tablas clínicas legibles en caja de 132 mm (170 mm - 20 mm - 18 mm).
\renewcommand{\arraystretch}{1.22}
\setlength{\tabcolsep}{4.5pt}
\AtBeginEnvironment{longtable}{\small\setlength{\LTpre}{0.55\baselineskip}\setlength{\LTpost}{0.55\baselineskip}}
\AtBeginEnvironment{tabular}{\small}

% Citas y bloques de evidencia o alerta.
\AtBeginEnvironment{quote}{%
  \small\color{libroprimario}%
}

% Ocultar el título automático de Quarto para usar portadilla propia.
\makeatletter
\renewcommand{\maketitle}{}
\makeatother

% Portada editorial vectorial integrada al PDF (170 × 244 mm = 6,69 × 9,61 in).
% Escala: 1 unidad = 0,1 mm (1700 × 2440 unidades).
\newcommand{\portadalibro}{%
  \begin{titlepage}%
  \thispagestyle{empty}%
  \begin{tikzpicture}[
    remember picture,
    overlay,
    shift={(current page.north west)},
    x=\paperwidth/1700,
    y=-\paperheight/2440
  ]
    % Fondo azul noche profundo (exactamente 170 × 244 mm).
    \fill[portadafondo] (0,0) rectangle (1700,2440);

    % Marco crema interior.
    \draw[portadapapel, opacity=0.22, line width=0.5pt] (96,96) rectangle (1604,2344);
    \draw[portadaoro, opacity=0.48, line width=0.9pt] (114,114) rectangle (1586,2326);

    % Banda superior de colección.
    \node[anchor=west, text=portadaoro, font=\sffamily\footnotesize\bfseries]
      at (164,205) {\librocoleccion};
    \draw[portadaoro, line width=0.8pt] (164,235) -- (920,235);

    % Título principal en tipografía de cartel, sin guiones.
    \node[anchor=north west, text=portadapapel, font=\displayfont\fontsize{28}{32}\selectfont\bfseries, text width=116mm, align=left]
      at (164,295) {\hyphenpenalty=10000 \exhyphenpenalty=10000 \librotitulo};

    % Subtítulo.
    \node[anchor=north west, text=portadapapel, opacity=0.90, font=\normalfont\fontsize{11}{15}\selectfont\itshape, text width=116mm, align=left]
      at (164,465) {\librosubtitulo};

    \draw[portadaoro, opacity=0.65, line width=0.7pt] (164,650) -- (1536,650);

    % Ilustración vectorial central: curva ROC, ejes bayesianos y nomograma de Fagan estilizado.
    \begin{scope}[shift={(850,1245)}]
      % Cuadrícula sutil de probabilidad.
      \draw[portadapapel, opacity=0.10, line width=0.5pt] (-380,-300) grid[step=95] (380,300);
      \draw[portadapapel, opacity=0.28, line width=0.8pt] (-380,300) -- (380,300);
      \draw[portadapapel, opacity=0.28, line width=0.8pt] (-380,-300) -- (-380,300);

      % Diagonal de azar (AUC = 0.5).
      \draw[portadapapel, opacity=0.25, dashed, line width=0.8pt] (-380,300) -- (380,-300);

      % Curva ROC de alta discriminación.
      \draw[portadacian, line width=2.0pt]
        (-380,300) .. controls (-360,-180) and (-160,-290) .. (380,-300);

      % Curva ROC moderada.
      \draw[portadaoro, opacity=0.75, line width=1.4pt]
        (-380,300) .. controls (-260,-40) and (-20,-220) .. (380,-300);

      % Tres ejes verticales del nomograma de Fagan superpuestos.
      \foreach \x/\lbl in {-260/{P_0}, 0/{\text{LR}}, 260/{P_1}} {
        \draw[portadaoro, opacity=0.45, line width=0.8pt] (\x,-230) -- (\x,230);
        \node[text=portadaoro, font=\sffamily\tiny\bfseries] at (\x,-260) {$\lbl$};
      }

      % Trayectoria pretest -> LR+ -> postest.
      \draw[portadaoro, line width=1.5pt] (-260,120) -- (0,-30) -- (260,-180);
      \foreach \px/\py in {-260/120, 0/-30, 260/-180} {
        \fill[portadaoro] (\px,\py) circle (9);
        \draw[portadapapel, line width=0.6pt] (\px,\py) circle (15);
      }

      % Etiquetas discretas.
      \node[text=portadaoro, opacity=0.90, font=\sffamily\tiny] at (-220,340) {SENSIBILIDAD · ESPECIFICIDAD};
      \node[text=portadaoro, opacity=0.90, font=\sffamily\tiny] at (220,340) {LR+ \enspace·\enspace LR− \enspace·\enspace FAGAN};
    \end{scope}

    % Tarjeta de la edición independiente.
    \fill[portadapapel, opacity=0.06] (164,1715) rectangle (1536,1935);
    \draw[portadaoro, opacity=0.45, line width=0.6pt] (164,1715) rectangle (1536,1935);
    \node[anchor=west, text=portadaoro, font=\sffamily\scriptsize\bfseries]
      at (205,1770) {PUBLICACIÓN INDEPENDIENTE · EDICIÓN CON DOI E ISBN PROPIOS};
    \node[anchor=west, text=portadapapel, opacity=0.92, font=\sffamily\scriptsize]
      at (205,1825) {4 capítulos de fundamentos \enspace·\enspace 48 condiciones verificadas de medsemiotics-db};
    \node[anchor=west, text=portadapapel, opacity=0.82, font=\sffamily\scriptsize]
      at (205,1880) {18 casos socráticos \enspace·\enspace 90 preguntas de autoevaluación \enspace·\enspace HTML, PDF y EPUB};

    % Pie de autoría y sello editorial.
    \node[anchor=north west, align=left, text=portadapapel, font=\sffamily\small\bfseries]
      at (164,2050) {\libroautoresportada};
    \node[anchor=west, text=portadaoro, font=\sffamily\scriptsize]
      at (164,2300) {\librosellopie};
  \end{tikzpicture}%
  \null\vfill
  \end{titlepage}%
  \cleardoublepage
}

% Entornos semánticos activados por el filtro Lua diseno-pdf.lua.
\newtcolorbox{triadasigno}{
  enhanced,
  breakable,
  colback=librofondo,
  colframe=libroprimario,
  boxrule=0.4pt,
  leftrule=2.6pt,
  arc=1.5pt,
  left=8pt,
  right=8pt,
  top=6pt,
  bottom=6pt,
  before skip=0.65\baselineskip,
  after skip=0.65\baselineskip,
  fontupper=\small
}

\newtcolorbox{alertaficha}{
  enhanced,
  breakable,
  colback=libroalerta,
  colframe=libroacento,
  boxrule=0.4pt,
  leftrule=2.8pt,
  arc=1.5pt,
  left=8pt,
  right=8pt,
  top=6pt,
  bottom=6pt,
  before skip=0.75\baselineskip,
  after skip=0.75\baselineskip,
  fontupper=\small
}

\newtcolorbox{casosocratico}{
  enhanced,
  breakable,
  colback=librosocratico,
  colframe=librosocraticoborde,
  boxrule=0.5pt,
  leftrule=3.0pt,
  arc=1.5pt,
  left=9pt,
  right=9pt,
  top=7pt,
  bottom=7pt,
  before skip=0.8\baselineskip,
  after skip=0.8\baselineskip,
  fontupper=\small
}

\newtcolorbox{perlasclinicas}{
  enhanced,
  breakable,
  colback=librofondo,
  colframe=libroprimario,
  boxrule=0.4pt,
  leftrule=2.6pt,
  arc=1.5pt,
  left=8pt,
  right=8pt,
  top=6pt,
  bottom=6pt,
  before skip=0.65\baselineskip,
  after skip=0.65\baselineskip,
  fontupper=\small
}

\newenvironment{portadillapdf}{%
  \thispagestyle{empty}%
  \vspace*{1.6cm}%
  \begin{flushleft}%
  \setlength{\parindent}{0pt}%
  \setlength{\parskip}{0.55\baselineskip}%
}{%
  \end{flushleft}%
  \vfill
  {\color{libroacento}\rule{0.42\linewidth}{0.8pt}\par}%
  \vspace{0.35\baselineskip}%
  {\sffamily\footnotesize\scshape\color{librogris}\librotitulo\par}%
  \clearpage
}

\newenvironment{portadillatitulo}{%
  \begingroup
  \displayfont\fontsize{25}{29}\selectfont\bfseries\color{libroprimario}%
  \hyphenpenalty=10000 \exhyphenpenalty=10000 \raggedright\noindent
}{%
  \par\endgroup
}

\newenvironment{creditospdf}{%
  \thispagestyle{empty}%
  \vspace*{\fill}%
  \begingroup
  \footnotesize\color{librogris}%
  \setlength{\parindent}{0pt}%
  \setlength{\parskip}{0.45\baselineskip}%
}{%
  \par\endgroup
  \clearpage
}

\newenvironment{enlaceficha}{%
  \par\vspace{0.25\baselineskip}%
  \noindent\begingroup\footnotesize\sffamily\color{librogris}%
}{%
  \par\endgroup\vspace{0.45\baselineskip}%
}
